%%%
% Cible 4 opérations
%%%
\def\filedateCibleQOp{2025/05/26}%
\def\fileversionCibleQOp{0.1a}%
\message{-- \filedateCibleQOp\space v\fileversionCibleQOp}%

\setKVdefault[CibleQO]{Couleur={},Rayon=2cm,ValeurMin=4,ValeurMax=10,Solution=false,DivE=false,Style="<->",Graine={},Cible={},Operations={}}%

\ExplSyntaxOn
\NewDocumentCommand\PfCCQODiviseurs{m}{%
  % #1 nombre donné
  % La liste \RetiensDiviseurs stocke les diviseurs
  % La liste \RetiensNonDiviseurs stocke les non diviseurs inférieurs à 10.
  \clist_clear:N \l_tmpa_clist
  \clist_clear:N \l_tmpb_clist
  \seq_clear:N \l_tmpa_seq
  \seq_clear:N \l_tmpb_seq
  \int_step_inline:nnn {2} {\fpeval{ceil(#1/2)}} {%
    \int_compare:nNnTF {\int_mod:nn {#1} {##1}} = {0} {%
      \seq_put_right:Ne \l_tmpa_seq {##1}
    } {\int_compare:nNnT {##1} < {10} {\seq_put_right:Ne \l_tmpb_seq {##1}}}
  }%
  \clist_set_from_seq:NN \l_tmpa_clist \l_tmpa_seq
  \clist_set_from_seq:NN \l_tmpb_clist \l_tmpb_seq
  \edef\RetiensDiviseurs{\clist_use:N \l_tmpa_clist}
  \edef\RetiensNonDiviseurs{\clist_use:N \l_tmpb_clist}
}
\ExplSyntaxOff

\newtoks\toklistecibleA%
\newtoks\toklistecibleab%
\newtoks\toklisteciblestyle%

\def\UpdatetoksCibleQO#1\nil{\addtotok\toklistecibleA{"#1",}}%
\def\UpdatetoksCibleQOO#1\nil{\addtotok\toklistecibleab{"#1",}}%
\def\UpdatetoksCibleStyle#1\nil{\addtotok\toklisteciblestyle{"#1"}}%%

%
\NewDocumentCommand\CibleQuatreOperations{o}{%
  \useKVdefault[CibleQO]%
  \setKV[CibleQO]{#1}%
  \ifemptyKV[CibleQO]{Graine}{}{\PfCGraineAlea{\useKV[CibleQO]{Graine}}}%
  \toklistecibleA{}%
  \toklistecibleab{}%
  \toklisteciblestyle{}%
  \expandafter\UpdatetoksCibleStyle\useKV[CibleQO]{Style}\nil%
  % Le style vaut :\the\toklisteciblestyle
  \ifemptyKV[CibleQO]{Couleur}{\colorlet{CouleurCibleQO}{Crimson}}{\setKV[CibleQO]{Solution}\colorlet{CouleurCibleQO}{\useKV[CibleQO]{Couleur}}}%
  % On définit le nombre cible aléatoirement ou fixement.
  \ifemptyKV[CibleQO]{Cible}{%
    \ChoixAlea{\useKV[CibleQO]{ValeurMin}}{\useKV[CibleQO]{ValeurMax}}{\PfCCiblePremierFacteur}%
    \ChoixAlea{\useKV[CibleQO]{ValeurMin}}{\useKV[CibleQO]{ValeurMax}}{\PfCCibleDeuxiemeFacteur}%
    \edef\PfCCibleQO{\fpeval{\PfCCiblePremierFacteur*\PfCCibleDeuxiemeFacteur}}%
  }{%
    \edef\PfCCibleQO{\useKV[CibleQO]{Cible}}%
  }%
  %  On définit les opérations : deux additions, deux soustractions, deux "divisions" inversées, deux "divisions" avec éventuellement une division euclidienne.
  \PfCCQODiviseurs{\PfCCibleQO}%
  \setsepchar{,}\ignoreemptyitems%
  \readlist*\PfCRetiensDiviseurs{\RetiensDiviseurs}%
  \reademptyitems%
  \ifemptyKV[CibleQO]{Operations}{%
    %% Il y a \PfCRetiensDiviseurslen{} diviseurs.
    \ifnum\PfCRetiensDiviseurslen=1\relax%
      \setKV[CibleQO]{DivE}%
    \fi%
    \ifboolKV[CibleQO]{DivE}{%
      \edef\PfCCQOListeOperations{+,+,-,-,*,*,/,d}%
    }{%
      \edef\PfCCQOListeOperations{+,+,-,-,*,*,/,/}%
    }%
    \edef\PfCCQOListeAddition{2,3,4,5,6,7,8,9,10,11,12,13,14,15,16}%
    \edef\PfCCQOListeMultiplication{2,3,4,5,6,7,8,9,10}%
    \edef\PfCFooNombres{}%
    \MelangeListe{\PfCCQOListeAddition}{4}%
    \edef\PfCFooNombres{\PfCFooNombres,\faa}%
    \MelangeListe{\PfCCQOListeMultiplication}{2}%
    \edef\PfCFooNombres{\PfCFooNombres,\faa}%
    \ifnum\PfCRetiensDiviseurslen=1\relax%carré parfait
      \edef\PfCFooNombres{\PfCFooNombres,\RetiensDiviseurs}%
      \MelangeListe{\RetiensNonDiviseurs}{1}%
      \edef\PfCFooNombres{\PfCFooNombres,\faa}%
    \else%
      \MelangeListe{\RetiensDiviseurs}{2}%
      \edef\PfCFooNombres{\PfCFooNombres,\faa}%
    \fi%
    %  % Les nombres choisis sont \PfCFooNombres\\
    \MelangeListe{1,2,3,4,5,6,7,8}{8}%
    \setsepchar{,}\ignoreemptyitems%
    \readlist*\PfCOrdreOperations{\faa}%
    \readlist*\PfCListeCibleNombres{\PfCFooNombres}%
    \readlist*\PfCListeCibleOperations{\PfCCQOListeOperations}%
    \reademptyitems%
    % On détermine les solutions
    \edef\PfCListeSolution{\fpeval{\PfCCibleQO-\PfCListeCibleNombres[1]},%
      \fpeval{\PfCCibleQO-\PfCListeCibleNombres[2]},%
      \fpeval{\PfCCibleQO+\PfCListeCibleNombres[3]},%
      \fpeval{\PfCCibleQO+\PfCListeCibleNombres[4]},%
      \fpeval{\PfCCibleQO*\PfCListeCibleNombres[5]},%
      \fpeval{\PfCCibleQO*\PfCListeCibleNombres[6]},%
      \fpeval{\PfCCibleQO/\PfCListeCibleNombres[7]}%
    }%
    \ifboolKV[CibleQO]{DivE}{%
      \edef\PfCListeSolution{\PfCListeSolution,\fpeval{floor(\PfCCibleQO/\PfCListeCibleNombres[8])}}%
    }{%
      \edef\PfCListeSolution{\PfCListeSolution,\fpeval{\PfCCibleQO/\PfCListeCibleNombres[8]}}%
    }%
    \setsepchar{,}\ignoreemptyitems%
    \readlist*\PfCListeCibleSolutions{\PfCListeSolution}%
    \reademptyitems%
    \edef\PfCFooRetiensTout{}%
    \edef\PfCFooRetiensToutAvecSol{}%
    \xintFor* ##1 in{\xintSeq{1}{8}}\do{%
      \edef\PfCCibleRang{\PfCOrdreOperations[##1]}%
      \edef\PfCCibleRang{\PfCCibleRang}%
      \StrCompare{\PfCListeCibleOperations[\PfCCibleRang]}{*}[\PfCTestEtoile]%
      \StrCompare{\PfCListeCibleOperations[\PfCCibleRang]}{/}[\PfCTestSlash]%
      \StrCompare{\PfCListeCibleOperations[\PfCCibleRang]}{d}[\PfCTestDE]%
      \xintifboolexpr{\PfCTestDE==0}{%
        \quotient{\PfCCibleQO}{\PfCListeCibleNombres[8]}%
        \modulo{\PfCCibleQO}{\PfCListeCibleNombres[8]}%
        \edef\PfCFooRetiensTout{\PfCFooRetiensTout,$\pointilles[5mm]\PfCSymbolTimes\num{\PfCListeCibleNombres[8]}\ifnum\remainder>0\relax+\num{\the\remainder}\fi$}%
        \edef\PfCFooRetiensToutAvecSol{\PfCFooRetiensToutAvecSol,$\mathcolor{CouleurCibleQO}{\num{\the\intquotient}}\PfCSymbolTimes\num{\PfCListeCibleNombres[8]}\ifnum\remainder>0\relax+\num{\the\remainder}\fi$}%
      }{%
        \edef\PfCFooRetiensTout{\PfCFooRetiensTout,$\pointilles[5mm]\xintifboolexpr{\PfCTestEtoile==0}{\PfCSymbolDiv}{\xintifboolexpr{\PfCTestSlash==0}{\PfCSymbolTimes}{\PfCListeCibleOperations[\PfCCibleRang]}}\num{\PfCListeCibleNombres[\PfCCibleRang]}$}%
        \edef\PfCFooRetiensToutAvecSol{\PfCFooRetiensToutAvecSol,$\mathcolor{CouleurCibleQO}{\num{\PfCListeCibleSolutions[\PfCCibleRang]}}\xintifboolexpr{\PfCTestEtoile==0}{\PfCSymbolDiv}{\xintifboolexpr{\PfCTestSlash==0}{\PfCSymbolTimes}{\PfCListeCibleOperations[\PfCCibleRang]}}\num{\PfCListeCibleNombres[\PfCCibleRang]}$}%
      }%
   }%
 }{% Les opérations sont données par l'utilisateur
   \edef\PfCOperationsListees{\useKV[CibleQO]{Operations}}%
   %  ListeOperations : \PfCOperationsListees\par
   \setsepchar{,}\ignoreemptyitems%
   \readlist*\PfCListeOperationsTestees{\PfCOperationsListees}%
   \reademptyitems%
   \ifnum\PfCListeOperationsTesteeslen<13\relax%
     \ifnum\PfCListeOperationsTesteeslen>9\relax\setKV[CibleQO]{Rayon=3cm}\fi%
     \StrCount{\PfCOperationsListees}{+}[\PfCNbAdd]%
     %Il y a \num{\PfCNbAdd}~additions\par
     \StrCount{\PfCOperationsListees}{-}[\PfCNbSous]%
     \StrCount{\PfCOperationsListees}{*}[\PfCNbMul]%
     \StrCount{\PfCOperationsListees}{/}[\PfCNbDiv]%
     \edef\PfCOperationsRetenues{}%
     \edef\PfCNombresRetenus{}%
     \edef\PfCSolutionsRetenues{}%
     \ifnum\PfCNbAdd>0\relax%
       \xintFor* ##1 in{\xintSeq{1}{\PfCNbAdd}}\do{%
         \edef\PfCOperationsRetenues{\PfCOperationsRetenues,+}%
       }%
       \PfCChoixAleaMultiMix{\PfCNbAdd}{1}{\fpeval{\PfCCibleQO-1}}{0}{\NbAddAv}%
       \edef\PfCNombresRetenus{\PfCNombresRetenus,\NbAddAv}%
     \fi%
     \ifnum\PfCNbSous>0\relax%
       \xintFor* ##1 in{\xintSeq{1}{\PfCNbSous}}\do{%
         \edef\PfCOperationsRetenues{\PfCOperationsRetenues,-}%
       }%
       \PfCChoixAleaMultiMix{\PfCNbSous}{1}{\fpeval{\PfCCibleQO-1}}{0}{\NbSousAv}%
       \edef\PfCNombresRetenus{\PfCNombresRetenus,\NbSousAv}%
     \fi%
     \ifnum\PfCNbMul>0\relax%
       \xintFor* ##1 in{\xintSeq{1}{\PfCNbMul}}\do{%
         \edef\PfCOperationsRetenues{\PfCOperationsRetenues,*}%
       }%
       \PfCChoixAleaMultiMix{\PfCNbMul}{2}{13}{0}{\NbMulAv}%
       \edef\PfCNombresRetenus{\PfCNombresRetenus,\NbMulAv}%
     \fi%
     \ifnum\PfCNbDiv>0\relax%
       \xintFor* ##1 in{\xintSeq{1}{\PfCNbMul}}\do{%
         \edef\PfCOperationsRetenues{\PfCOperationsRetenues,/}%
       }%
       \PfCChoixAleaMultiMix{\PfCNbMul}{2}{13}{0}{\NbDivAv}%
       \edef\PfCNombresRetenus{\PfCNombresRetenus,\NbDivAv}%
     \fi%
     \setsepchar{,}\ignoreemptyitems%
     \readlist*\ListeOperationsRetenuesVraiment{\PfCOperationsRetenues}%
     \readlist*\ListeNombresRetenusVraiment{\PfCNombresRetenus}%
     \reademptyitems%
     \edef\PfCFooRetiensTout{}%
     \edef\PfCFooRetiensToutAvecSol{}%
     \foreachitem\comptenum\in\ListeOperationsRetenuesVraiment{%
       \IfStrEqCase{\comptenum}{%
         {+}{%
           \edef\PfCFooRetiensTout{\PfCFooRetiensTout,$\pointilles[5mm]+\ListeNombresRetenusVraiment[\comptenumcnt]$}%
           \edef\PfCFooRetiensToutAvecSol{\PfCFooRetiensToutAvecSol,$\mathcolor{CouleurCibleQO}{\num{\fpeval{\PfCCibleQO-\ListeNombresRetenusVraiment[\comptenumcnt]}}}+\ListeNombresRetenusVraiment[\comptenumcnt]$}%
         }%
         {-}{%
           \edef\PfCFooRetiensTout{\PfCFooRetiensTout,$\pointilles[5mm]-\ListeNombresRetenusVraiment[\comptenumcnt]$}%
           \edef\PfCFooRetiensToutAvecSol{\PfCFooRetiensToutAvecSol,$\mathcolor{CouleurCibleQO}{\num{\fpeval{\PfCCibleQO+\ListeNombresRetenusVraiment[\comptenumcnt]}}}-\ListeNombresRetenusVraiment[\comptenumcnt]$}%
         }%
         {*}{%
           \edef\PfCFooRetiensTout{\PfCFooRetiensTout,$\pointilles[5mm]\PfCSymbolDiv\ListeNombresRetenusVraiment[\comptenumcnt]$}%
           \edef\PfCFooRetiensToutAvecSol{\PfCFooRetiensToutAvecSol,$\mathcolor{CouleurCibleQO}{\num{\fpeval{\PfCCibleQO*\ListeNombresRetenusVraiment[\comptenumcnt]}}}\PfCSymbolDiv\ListeNombresRetenusVraiment[\comptenumcnt]$}%
         }%
         {/}{%
           \quotient{\PfCCibleQO}{\ListeNombresRetenusVraiment[\comptenumcnt]}%
           \modulo{\PfCCibleQO}{\ListeNombresRetenusVraiment[\comptenumcnt]}%
           \ifnum\remainder=0\relax\edef\PfCFooRetiensTout{\PfCFooRetiensTout,$\pointilles[5mm]\times\ListeNombresRetenusVraiment[\comptenumcnt]$}\else\edef\PfCFooRetiensTout{\PfCFooRetiensTout,$\pointilles[5mm]\PfCSymbolTimes\ListeNombresRetenusVraiment[\comptenumcnt]+\the\remainder$}\fi%
           \ifnum\remainder=0\relax\edef\PfCFooRetiensToutAvecSol{\PfCFooRetiensToutAvecSol,$\mathcolor{CouleurCibleQO}{\num{\the\intquotient}}\times\ListeNombresRetenusVraiment[\comptenumcnt]$}\else\edef\PfCFooRetiensToutAvecSol{\PfCFooRetiensToutAvecSol,$\mathcolor{CouleurCibleQO}{\num{\the\intquotient}}\PfCSymbolTimes\ListeNombresRetenusVraiment[\comptenumcnt]+\the\remainder$}\fi%
         }%
       }%
     }%
   \fi%
   \MelangeListe{\PfCFooRetiensTout}{\PfCListeOperationsTesteeslen}%
   \edef\PfCFooRetiensTout{\faa}%
   \MelangeListe{\PfCFooRetiensToutAvecSol}{\PfCListeOperationsTesteeslen}%
   \edef\PfCFooRetiensToutAvecSol{\faa}%
 }%
 \setsepchar{,}\ignoreemptyitems%
 \readlist*\PfCCibleQOListeEnonce{\PfCFooRetiensTout}%
 \readlist*\PfCCibleQOListeSolution{\PfCFooRetiensToutAvecSol}%
 \reademptyitems%
 \foreachitem\compteur\in\PfCCibleQOListeEnonce{\expandafter\UpdatetoksCibleQO\compteur\nil}%
 \foreachitem\compteur\in\PfCCibleQOListeSolution{\expandafter\UpdatetoksCibleQOO\compteur\nil}%
 \BuildCibleQuatreOperations{\PfCCibleQO}{\the\toklistecibleA}{\the\toklistecibleab}{\useKV[CibleQO]{Style}}%
}%

\def\BuildCibleQOCode{%
  vardef LectureOperations(text t)=
    nbop:=0;
    for p_=t:
      Op[nbop]=image(
        label(TEX(p_),(0,0));
      );
      nbop:=nbop+1;
    endfor;
  enddef;
}%

\NewDocumentCommand\BuildCibleQuatreOperations{mmmm}{%
  \mplibforcehmode%
  \begin{mplibcode}%
    \BuildCibleQOCode
    string Style;
    Style=#4;
    
    Cible=#1;
    
    picture Op[];
    
    Rayon:=\useKV[CibleQO]{Rayon};
    boolean Solution;
    Solution=\useKV[CibleQO]{Solution};
    
    if Solution:
      LectureOperations(#3)
    else:
      LectureOperations(#2)
    fi;
      
    path cc;
    pair O,A[];
    O=(0,0);
    picture CIBLE;
    CIBLE=image(label(TEX("\Huge\bfseries"&decimal(Cible)),O));
    draw CIBLE;
    cc=cercles(O,Rayon);
    path Fleche;
    for k=0 upto nbop-1:
      if (k*(360/nbop) mod 180)=0:
        A[k]=1.2[O,pointarc(cc,k*(360/nbop))];
      else:
        A[k]=1.1[O,pointarc(cc,k*(360/nbop))];
      fi;
      Fleche:=((0,0)--(center Op[k] shifted A[k])) cutbefore cercles(O,0.45*Rayon) cutafter cercles(O,0.90*Rayon);
      if Style="<->": 
        drawdblarrow (point(0.75*length Fleche) of Fleche)--(point(0) of Fleche);
      elseif Style="<-":
        drawarrow (point(0.75*length Fleche) of Fleche)--(point(0) of Fleche);
      elseif Style="->":
        drawarrow (point(0) of Fleche)--(point(0.75*length Fleche) of Fleche);
      elseif Style="-":
        draw (point(0) of Fleche)--(point(0.75*length Fleche) of Fleche);
      elseif Style="--":
        label(TEX("$=$") rotated (k*45),point(0.5*length Fleche) of Fleche);
      fi;
      trace Op[k] shifted (A[k]-center Op[k]);
    endfor;
  \end{mplibcode}%
}%
